\chapter{第二十讲\quad 危害公共安全罪}
\label{lecture:20}
\noindent{\small\hyperref[ref:lecture:20]{查阅本讲扩展规则}}\par

令 $a$ 为行为，$d$ 为危险，$r$ 为结果。$P(a,d)$ 表示行为产生指向公共安全的危险，$R(a,r)$ 表示依法达到严重程度且能归责于该行为的实际结果。教材的公共性检验看多数人或不特定人的生命、健康、重大公私财产是否受威胁，尤其看危险的传播范围、控制可能和现实暴露。人数是一项证据。题目已经给出具体罪名的专门条件时，仍须按相应条款逐一验证。

\section{第一节\quad 危险方法型犯罪}

\paragraph{法条与模型}
第114、115条的结构可写成
\[
\begin{aligned}
 \text{故意}\land\text{法定方法}\land P\land\neg R
    &\Longrightarrow \text{第114条路径},\\
 \text{故意}\land\text{法定方法}\land P\land R
    &\Longrightarrow \text{第115条第一款路径}.
\end{aligned}
\]
$\neg R$ 在此表示已经确认没有可归责的严重结果。严重结果尚不能证明时，不能把证据缺口当作 $\neg R$ 的事实证明。法定方法包括放火、决水、爆炸、投放危险物质，以及与列举行为危险性相当的其他危险方法。投放一项同时核对物质性质与传播路径。第115条第二款的过失路径另证罪过和实际严重后果。教材的放火既遂标准中，独立燃烧是可用的危险证据，不能省略公共性。已着手而危险未出现的情形，另审查未遂。

\paragraph{推导与反例}
甲故意点燃有人居住的相连建筑内的可燃物。鉴定确认火势已蔓延并现实危及多个住户，及时扑灭，未造成严重后果。由蔓延与住户暴露得 $P$，由结果证据得 $\neg R$，由故意放火和第114条推得该路径。若同样点火仅发生在与外界隔绝的环境，证据证明危险只针对特定物品且没有公共扩散可能，则 $P$ 缺失，需另查毁坏财物、伤害或杀人等。若改用不具危险性质的粉末引起恐慌，不可只凭「害怕」推导投放危险物质，应转查其他可能的规定。

\section{第二节\quad 破坏型犯罪}

\paragraph{法条与模型}
第116、117条同查足以造成交通工具倾覆、毁坏的危险 $D$，但对象不同。$V$ 表示作为直接破坏对象的法定交通工具，$F$ 表示交通设施或条文涵盖的其他破坏活动。分别以 $V\land D$、$F\land D$ 审查两条。第118条查电力、燃气等设备和公共危险。造成严重后果的转查第119条；过失的，另证第二款所需严重后果和过失。教材中「正在使用」要看功能状态、是否已投入使用等事实，不能机械等同于行为当刻正在移动。

\paragraph{推导与反例}
甲故意破坏次日按计划通车且未被隔离的铁路设施。题设鉴定足以造成列车倾覆，运行前及时发现，没有严重后果。设施给出 $F$，鉴定与运行计划给出 $D$，从而进入第117条，无须先有列车实际倾覆。若直接破坏的是运行中的列车并达到同一危险程度，则改变的是对象谓词 $V$，改查第116条。若对象是彻底退役、已脱离交通用途的零件，且破坏并不产生相关危险，单靠「曾属列车」不能推出 $D$。拆取电力设备还可能同时触发财产犯罪，应先查使用状态与公共危险，再作罪数评价。

\section{第三节\quad 交通型犯罪}

\paragraph{交通肇事的三层}
第133条与法释〔2000〕33号把判断分成基础、升格和逃逸致死三层。基础层 $B$ 检查公共交通管理范围、违反规范、符合司法解释的事故责任与结果组合，以及可归责性。升格层 $F$ 分两条路径。普通逃逸须有解释第2条第一款或第二款第（一）至（五）项所列基础情形，并明知肇事而为逃避追究逃跑。其他特别恶劣情节按第4条查证。若同一逃逸事实已作为第2条第二款第（六）项的入罪条件，不能再将该事实直接用作第3条的升格条件。死亡层 $E$ 要证明逃逸使被害人得不到救助而死亡，并审查死亡结果的罪过。
\[
 E:\quad
 \text{为逃避追究而逃逸}
 \leadsto\text{未获及时救助}
 \leadsto\text{死亡}.
\]
应说明报告、救助等行为的可行性和结果避免可能。对于适用逃逸致死档时，逃逸前的事故本身是否必须独立达到基本犯门槛，教材有争议讨论，应注明采用的解释，不能以逃逸致死反推先前事故门槛已满足。把伤者带离后隐藏、遗弃，致其无法获救而死亡或严重残疾，另按该解释第6条及相应罪过事实审查故意杀人、故意伤害。解释第5条第二款针对单位主管人员、车辆所有人、承包人或乘车人指使逃逸并致被害人因得不到救助而死亡规定特殊共犯处理。普通共同过失则应对各行为人的注意义务及因果作用分别评价。最高法指导性案例269号提示，不能仅凭事后逃逸导致的事故责任认定推定其对事故发生的刑法因果责任。

\paragraph{逐层推导}
甲违反交通运输管理法规造成一死一重伤，依法负事故主要责任，事故与其违规有归责联系，故先证得 $B$。甲知悉事故，为逃避追究离开，证得 $F$。若重伤者当时尚能救活，甲本可报警，却因逃逸无人及时救助并因此死亡，且对该死亡具备相应过失，则可依题设证明 $E$。改动一个事实：鉴定明确伤者在甲逃逸前已死亡，则 $\text{未获及时救助}\leadsto\text{死亡}$ 不成立。$E$ 失效，但不影响原先已经证明的 $B$，$F$ 仍按前提事实单独判断。

\paragraph{危险驾驶与妨害安全驾驶}
第133条之一先查道路、机动车和驾驶，再在四支中择一：情节恶劣的追逐竞驶、醉酒驾驶、校车业务或旅客运输的严重超员超速、违规运输危险化学品且危及公共安全。第三、四项的机动车所有人、管理人另查直接责任。2023年醉驾办案意见对道路公共性、酒精鉴定和情节显著轻微等有规定。因此，在向社会机动车开放的道路上醉酒驾驶且题设排除该等例外，可成立危险驾驶，未造成人员伤亡不单独阻断该支。仅在不向社会开放的封闭内部场所移车，则先复核道路条件。第133条之二的乘客路径检查行驶中的公共交通工具、暴力或抢控操纵装置、正常行驶受干扰及公共危险。驾驶人员擅离职守殴打他人或互殴有单独路径。若抢控使载客车辆偏离车道并危及公共安全，依相应要件推导该罪候选；若仅在停运车辆内发生争吵，则行驶和危险等条件不能沿用。两条均有同时触犯其他犯罪时择重的专门规定。

\section{第四节\quad 恐怖型犯罪}

\paragraph{法条与模型}
令 $O$ 为已查实的恐怖活动组织，$J$ 为组织、领导或参加，$H$ 为资助或招募、运送人员等本条评价的帮助，$Q$ 为第120条之二列举的准备行为。依第120条、第120条之一、第120条之二分开判断：
\[
 \begin{array}{ccl}
 O\land J&\Longrightarrow&\text{组织、领导、参加路径},\\
 \text{法定帮助对象}\land H&\Longrightarrow&\text{独立帮助路径},\\
 \text{相应目的与行为}\land Q&\Longrightarrow&\text{独立准备路径}.
 \end{array}
\]
把帮助和准备作为独立实行行为，是教材所称正犯化。前两项不以实际恐怖活动发生为一律必要条件。对组织、领导、参加后又实施杀人、爆炸、绑架等，第120条第二款有并罚规定。准备行为同时构成其他犯罪时，第120条之二第二款规定择重。教唆、帮助这些独立犯罪，另按总则检查其共犯条件。

\paragraph{推导与反例}
甲明知是恐怖活动组织，向该组织提供经查实的物质资助且对方已接收，尚未发生计划中的活动。组织性质、明知和资助交付共同满足 $H$ 路径，不能从「活动未发生」推出「独立帮助未完成」。若甲只与人员一般交往，组织关联、知情和资助等事实均未证实，则不能把交往自动改写成 $H$。一般阅读资料也不能直接证明第120条之二的特定准备行为与目的。

\section{第五节\quad 枪支类犯罪}

\paragraph{法条与模型}
先证 $G$（依鉴定及适用规则符合枪支等法定对象），再证来源、权限及行为类型。第125条规定非法制造、买卖、运输、邮寄、储存枪支、弹药、爆炸物等。该条第二款对毒害性、放射性、传染病病原体等物质没有列「邮寄」，另要求危害公共安全。第127条针对盗窃、抢夺、抢劫等取得方式，并分对象和条款。第128条第一款针对违反管理规定的非法持有、私藏。教材把无合法资格而擅自持有与原合法配备、配置者在条件消失后拒不交出分别分析，仍须核实解释中的对象、数量和证据标准。出租、出借的两个分支是
\[
\begin{aligned}
 \text{依法配备公务用枪}\land\text{非法出租或出借}
   &\Longrightarrow \text{第128条第二款},\\
 \text{依法配置枪支}\land\text{非法出租或出借}\land\text{严重后果}
   &\Longrightarrow \text{第128条第三款}.
\end{aligned}
\]
第129条另查依法配备公务用枪、丢失、不及时报告和严重后果及其归责联系。

\paragraph{推导与反例}
甲依法配备公务用枪，将经鉴定的公务枪支非法出借，尚未造成严重后果。身份与枪支性质建立第一支左项，出借建立右项，故按第二款审查。若只改成依法配置的民用枪支，严重后果仍未发生，则第三款的最后一项缺失。枪状外观也不能独立推出 $G$。若甲通过盗窃取得后继续持有，先分别辨认取得与持有的事实，再判断持有是否已被前罪评价，不能机械把两项时间状态计为两罪。

\section{第六节\quad 安全生产领域犯罪}

\paragraph{法条与模型}
设 $V$ 是生产、作业中违反安全管理规定，$L$ 是第134条之一列举的三类情形之一，$D$ 是发生重大伤亡或其他严重后果的现实危险，$R$ 是已经发生的法定严重后果且能归责。比较
\[
\begin{aligned}
 V\land L\land D&\Longrightarrow \text{危险作业入口（第134条之一）},\\
 V\land R&\Longrightarrow \text{重大责任事故入口（第134条第一款）},\\
 \text{安全设施或条件不符}\land R\land\text{直接责任}
   &\Longrightarrow \text{重大劳动安全事故入口（第135条）}.
\end{aligned}
\]
第134条第二款另查强令违章冒险作业，或者明知重大事故隐患而不排除、仍冒险组织作业，及相应严重后果。危险与实害分别记录，并按各罪要求核验罪过，不能只因结果发生就把危险罪机械替换为事故罪。第139条之一在事故之后检查负有报告职责、不报或谎报、因此贻误抢救且情节严重。

\paragraph{推导与反例}
甲在有重大事故隐患的作业中被依法责令停止施工，却明知而拒不执行。证据证明作业持续已带来现实的重大伤亡危险，尚无实际事故。由违法作业得 $V$，由责令和拒绝得 $L$ 的第二支，由险情证据得 $D$，得危险作业路径。若只是一般违章，既未满足列举情形也无现实危险，不能单凭 $V$ 推出该罪。若事故后负责报告的人谎报，但题设证明全部抢救早已结束且没有贻误，则第139条之一的因果条件缺失，仍可另查其他责任。

\section*{闭卷记忆与验算}
\paragraph{复述路线}
写六列「方法、对象、领域、危险、结果、归责」。先用方法和对象匹配条文，再查危险究竟指向谁以及如何扩散。把「有危险」「有实际结果」「结果可归责」分别落格。最后查特殊主体、独立化行为、升格与择重。速记「先找对象，再问危险，实害另证，逃逸补链」。「逃逸补链」要求逐箭头写出逃逸、未获救助和死亡。

\begin{enumerate}
\item 题：甲破坏运行中的交通工具，鉴定足以倾覆，但车在事故前停下。为何仍可成立犯罪？答：按第116条证明对象和足以倾覆、毁坏的危险即可进入基本路径。尚无严重后果只排除相应的第119条结果入口。
\item 题：肇事者逃离现场，医学上即使立即救助也不能避免死亡，怎样处理？答：将逃逸、逃逸与缺乏救助、缺乏救助与死亡分别证明。最后一段不能成立，不能适用因逃逸致人死亡档。基础罪和肇事后逃逸档另查。
\item 题：生产作业中只证实一项安全违规，能否直接推出危险作业罪？答：不能。还要指明第134条之一的列举情形和重大伤亡或其他严重后果的现实危险，再验证相应罪过。
\end{enumerate}
